\section{Conclusion}
Whether quantization undermines concept erasure depends on the bit width, and the
reason can be computed before generating any images. For closed-form cross-attention
edits, image behavior is a function of the key/value outputs, and the set of
checkpoints that deploy exactly as the original is a box. Together these turn the
reconstruction question into a certifiable convex program over the original-pattern
box. At 8 bits that box is narrow: the best reconstruction reproduces about 40\% of the
erasure's key/value effect, and the resulting checkpoint scores like the unedited model
in FP16. At 4 bits it is wide enough that the constructed checkpoint scores close to
the honest erasure in FP16 yet restores the erased concept after deployment. Honest erasures show no significant deployment gap at any tested
precision, but whole-UNet INT4 hints at a prompt-dependent passive re-emergence of
nudity that deserves a powered follow-up. Deployers who quantize
erased models to 4 bits should audit the quantized model, or run the image-free
key/value audit against the base model, rather than trust a full-precision
evaluation.

\section*{Open Science}
All code, per-image scores, run manifests with result and source hashes, the ordered
prompt lists, and the scripts that regenerate every table and figure from the stored
outputs are provided in the anonymous artifact at \url{https://anonymous.4open.science/r/SaTML-7863/}. The CPU
feasibility analysis (Section~\ref{sec:feasibility}) reruns from the public SD-1.5
weights and the stored calibration statistics without a GPU. Base models, detectors,
and prompt sets are obtained from their original distributions under their licenses;
generated images are not redistributed.

\section*{LLM Usage Considerations}
An LLM-based coding assistant helped with source inspection, analysis scripts, and
editing of this manuscript. All experiments were executed by the authors'
code, all numbers are produced by deterministic scripts from stored outputs, and the
authors verified the methods, proofs, references, and conclusions.

\section*{Ethical Considerations}
The nudity study measures a detector-defined suppression target and does not equate
nudity with harm. The pipeline's safety checker is disabled to measure the edited
generator itself. Generated images are scored in temporary storage, deleted, and not
released. Our attack construction is a known class (quantization-conditioned
weights~\cite{egashira2024}) applied to erasure. We release it together with the
certificate and the image-free audit of Section~\ref{sec:audit}, which lets a deployer
flag such checkpoints before deployment without generating images.
